\subsubsection{2.2.4 Symmetrie, Symmetriebruch und die Rolle der Projektion}

Nachdem sich in den numerischen Untersuchungen zunehmend herauskristallisierte, dass die beobachtete Raumzeit wahrscheinlich nicht direkt dem OPP entspricht, sondern aus einer geeigneten Projektion hervorgeht, stellte sich eine neue Frage:

\begin{center}
\emph{Kann man die Projektion nicht von der bekannten Physik aus rückwärts rekonstruieren?}
\end{center}

Bisher lautete der Ansatz:

\[
\mathcal{O}
\longrightarrow
\Pi
\longrightarrow
\mathcal{M}
\]

also

\begin{center}
``Aus dem OPP die Physik erzeugen.''
\end{center}

Die neue Idee kehrt diese Richtung teilweise um.

Da unsere Physik experimentell hervorragend bestätigt ist, kennen wir bereits zahlreiche Eigenschaften des Zielraumes.

Insbesondere kennen wir seine Symmetrien.

Damit entsteht die Fragestellung

\[
\boxed{
\text{Welche Eigenschaften muss eine Projektion besitzen,
damit genau diese Symmetrien emergieren?}
}
\]

Dies markiert einen Übergang vom rein konstruktiven Ansatz zu einem rekonstruktiven Ansatz.

\section*{Symmetrien der beobachteten Physik}

Die heutige Physik besitzt eine bemerkenswerte Vielzahl hochentwickelter Symmetrien.

Unter anderem

\[
SO(1,3)
\]

als Lorentzsymmetrie der Raumzeit,

sowie die Eichgruppen

\[
SU(3),
\qquad
SU(2),
\qquad
U(1).
\]

Diese Gruppen treten nicht zufällig auf.

Sie bestimmen unmittelbar

\begin{itemize}

\item die Erhaltungssätze,

\item die Wechselwirkungen,

\item die Elementarteilchen,

\item die Eichfelder,

\item die Dynamik des Standardmodells.

\end{itemize}

Daraus ergibt sich eine naheliegende Überlegung.

Falls unsere Physik tatsächlich eine Projektion des OPP darstellt,

dann muss die Projektion diese Symmetrien entweder

\begin{enumerate}

\item erhalten,

oder

\item erzeugen.

\end{enumerate}

\section*{Der OPP als hochsymmetrischer Raum}

Die numerischen Untersuchungen sprechen dafür,

dass der OPP zunächst eine wesentlich höhere effektive Dimension besitzt,

typischerweise

\[
d_{\mathrm{eff}}
\approx
30
-
45.
\]

Ein derart hochdimensionales System besitzt zwangsläufig erheblich größere Freiheitsgrade als unsere beobachtete Raumzeit.

Dies legt die Vermutung nahe,

dass der OPP selbst eine sehr große ursprüngliche Symmetrie besitzt.

Symbolisch

\[
G_{\mathrm{OPP}}
\supset
SO(1,3)
\times
SU(3)
\times
SU(2)
\times
U(1).
\]

Dabei ist

\[
G_{\mathrm{OPP}}
\]

keine konkret bekannte Gruppe,

sondern lediglich eine Platzhalterbezeichnung für die fundamentale Symmetrie des OPP.

\section*{Symmetriebruch}

Die beobachtete Physik besitzt offensichtlich nicht die vollständige Symmetrie des OPP.

Folglich muss zwischen beiden Ebenen ein Symmetriebruch stattfinden.

Formal

\[
G_{\mathrm{OPP}}
\longrightarrow
G_{\mathrm{Physik}}.
\]

Dieser Übergang erinnert unmittelbar an spontane Symmetriebrechung.

Die zentrale Frage lautet jedoch:

\begin{center}

\emph{Wo findet dieser Symmetriebruch statt?}

\end{center}

Die PST schlägt eine neue Antwort vor.

Nicht der OPP selbst bricht seine Symmetrie.

Sondern

\[
\boxed{
\text{die Projektion bricht die Symmetrie.}
}
\]

Der OPP bleibt vollständig erhalten.

Lediglich die beobachtbare Projektion besitzt eine reduzierte Symmetrie.

\section*{Projektionsinduzierter Symmetriebruch}

Sei

\[
\Pi
\]

eine Projektion.

Dann besitzt allgemein

\[
\Pi(G_{\mathrm{OPP}})
\neq
G_{\mathrm{OPP}}.
\]

Es bleibt lediglich eine Untergruppe erhalten.

Formal

\[
G_{\Pi}
=
\Pi(G_{\mathrm{OPP}})
\subseteq
G_{\mathrm{OPP}}.
\]

Diese Untergruppe wäre genau die Symmetrie,

die wir experimentell beobachten.

Damit verschiebt sich die Frage

von

\begin{quote}

``Warum besitzt die Natur diese Symmetrien?''

\end{quote}

zu

\begin{quote}

``Warum erhält gerade diese Projektion genau diese Untergruppen?''

\end{quote}

\section*{Innere Freiheitsgrade}

Ein weiterer Gedanke entstand aus den numerischen Ergebnissen.

Die Diffusionsprojektion erzeugte eine effektive Vierdimensionalität,

ohne den hochdimensionalen OPP vollständig zu eliminieren.

Damit bleiben zwangsläufig Freiheitsgrade übrig.

Diese könnten physikalisch interpretiert werden als

\begin{itemize}

\item innere Freiheitsgrade,

\item Eichfreiheitsgrade,

\item interne Symmetrien,

\item Quantenzahlen.

\end{itemize}

Die zusätzlichen Dimensionen verschwinden also möglicherweise gar nicht.

Sie werden lediglich

\begin{center}

\emph{für den makroskopischen Beobachter unsichtbar.}

\end{center}

Dies erinnert an Kaluza-Klein-Theorien,

ohne deren geometrische Voraussetzungen übernehmen zu müssen.

\section*{Die Rolle komplexer Relationen}

Während der Diskussion entstand außerdem eine weitere Vermutung.

Die bisherigen Simulationen verwenden ausschließlich reelle Relationen

\[
R_{ij}
\in
\mathbb{R}.
\]

Viele fundamentale Symmetrien der Physik beruhen jedoch auf komplexen Gruppen,

insbesondere

\[
SU(n).
\]

Dies führte zur Hypothese,

dass zukünftige Versionen des OPP möglicherweise komplexwertige Relationen verwenden sollten,

also

\[
R_{ij}
\in
\mathbb{C}.
\]

Dadurch würden zusätzlich Phaseninformationen entstehen.

Diese könnten die Grundlage für

\begin{itemize}

\item Interferenz,

\item Quantenphasen,

\item Eichfelder,

\item unitäre Symmetrien

\end{itemize}

bilden.

Ob dies tatsächlich erforderlich ist,

bleibt eine offene Forschungsfrage.

\section*{Top-Down statt Bottom-Up}

Die gesamte Diskussion führte schließlich zu einem methodischen Wechsel.

Der ursprüngliche Forschungsansatz lautete

\begin{center}

``Wir erzeugen beliebige Projektionen und hoffen,

dass daraus irgendwann Physik entsteht.''

\end{center}

Der neue Ansatz lautet

\begin{center}

``Wir kennen bereits die Physik.

Nun suchen wir genau diejenigen Projektionen,

welche diese Physik reproduzieren.''

\end{center}

Damit wird die bekannte Physik selbst zu einer experimentellen Randbedingung.

Nicht jede mathematisch mögliche Projektion ist zulässig.

Zulässig sind ausschließlich diejenigen,

welche

\begin{itemize}

\item Lorentz-Invarianz,

\item stabile Raumzeit,

\item bekannte Eichsymmetrien,

\item beobachtete Quantisierung,

\item die bekannten Erhaltungssätze

\end{itemize}

gleichzeitig reproduzieren.

\section*{Ausblick}

Die Symmetriebetrachtungen liefern noch keine vollständige Theorie.

Sie verändern jedoch den Blickwinkel grundlegend.

Die Projektion ist nicht länger nur ein mathematisches Werkzeug.

Sie wird selbst zu einem physikalischen Objekt.

Damit ergibt sich als langfristiges Ziel:

\[
\boxed{
\text{Nicht der OPP allein bestimmt die Physik,
sondern die Kombination aus OPP und zulässiger Projektion.}
}
\]

Die numerischen Simulationen der folgenden Kapitel dienen genau dazu,

diese zulässigen Projektionen schrittweise einzugrenzen.

